\documentclass[11pt]{article}
\usepackage[a4paper,margin=28mm]{geometry}
\usepackage{amsmath,amssymb,amsthm,booktabs,hyperref}
\newtheorem{theorem}{Theorem}[section]
\newtheorem{lemma}[theorem]{Lemma}
\newtheorem{proposition}[theorem]{Proposition}
\newtheorem{corollary}[theorem]{Corollary}
\newcommand{\dist}{\operatorname{dist}}
\title{Regular graphs of girth four and diameter divisible by four}
\author{Human Compression project: mathematical preparation}
\date{31 October 2026}
\begin{document}\maketitle

\paragraph{Scope and attribution.}
The statements in this note are proved by hand from a construction and
an equality argument proposed by the project's coordinating agent.
The independent-set degree bound, geodesic shortcuts and separator
distance arguments are elementary standard methods. The contribution
of this note is their explicit common assembly, including all degree,
girth and diameter checks for the proposed constructions. No assertion
of literature novelty is made. No original numerical census is used
as a premise or stopping condition. No Source, graph enumeration,
scientific data, Lean proof or TeX build was executed. The note awaits
independent mathematical review and adoption.

All graphs are finite, simple and connected. For integers $k\geq3$ and
$d\geq1$, let $n(k;4,d)$ denote the minimum order of a $k$-regular graph
of girth exactly four and diameter exactly $d$, when this domain is
nonempty. The lower bounds below hold for every triangle-free graph,
without requiring a four-cycle. The explicit upper constructions do
have four-cycles.

\begin{theorem}\label{main}
For every $k\geq3$ and $r\geq2$,
\[
 n(k;4,4r)=
 \begin{cases}
 2k(r+1),&k\text{ composite},\\
 2k(r+1)+2,&k\text{ prime}.
 \end{cases}
\]
\end{theorem}

Degree two is deliberately excluded. A connected simple 2-regular
graph is a cycle; it cannot have girth four and diameter $4r$ for
$r\geq2$. If only triangle-freeness is required, the cycle of diameter
$4r$ has order $8r$, a different formula.

\section{The independent layers and the equality obstruction}

Fix a geodesic $p_0,p_1,\ldots,p_{4r}$ in a triangle-free $k$-regular
graph $\Gamma$. For $0\leq i\leq r$ write
\[
 z_i=p_{4i},\qquad I_i=N(z_i),\qquad I=\bigcup_{i=0}^r I_i.
\]

\begin{lemma}\label{lower}
$I$ is an independent set of size $(r+1)k$. Consequently
$|V\Gamma|\geq2k(r+1)$. If equality holds, its complement $J$ is also
independent, and $I,J$ form the two equal parts of a bipartite graph.
\end{lemma}
\begin{proof}
Each $I_i$ is independent by triangle-freeness. Distinct $I_i,I_j$
are disjoint, because a common vertex gives a length-two path between
$z_i,z_j$, whose distance is $4|i-j|\geq4$. An edge between them
would give a length-three path. Thus their union is independent and
has the stated size. All $k|I|$ edges incident with it go into the
complement, whose total degree is $k(|V\Gamma|-|I|)$. This proves
the bound. At equality all complement degrees are used by these
crossing edges, so there are no complement edges.
\end{proof}

Assume equality for the rest of this section. A vertex of $J$ can
have neighbors in at most two of the sets $I_i$, and if it meets two,
they must be consecutive. Indeed if it met $I_i$ and $I_j$ with
$|i-j|\geq2$, there would be a path of length four between $z_i,z_j$,
shorter than their geodesic distance of at least eight. Thus every
vertex of $J$ has one of the following types:
\[
 X_i:\ N(x)\subseteq I_i,\qquad
 F_i:\ N(x)\cap I_i\ne\varnothing\ne N(x)\cap I_{i+1}.
\]
There are no other types. Every member of $X_i$ is adjacent to all
$k$ vertices of $I_i$, since its degree is $k$. Thus $X_i$ is a
false-twin class and contains $z_i$. Write
\[
 x_i=|X_i|\geq1,\qquad q_i=k-x_i.
\]

\begin{lemma}\label{interfaces}
For each $i$, $1\leq q_i\leq k-1$. For $1\leq i\leq r-1$, the set
$I_i$ splits into nonempty sets $I_i^-,I_i^+$, of vertices whose
neighbors outside $X_i$ lie respectively in $F_{i-1},F_i$.
All of $I_0$ uses $F_0$, and all of $I_r$ uses $F_{r-1}$.
\end{lemma}
\begin{proof}
Every vertex of $I_i$ is adjacent to all of $X_i$ and therefore has
exactly $q_i$ other neighbors. If $x_i=k$, every vertex of $I_i$ is
saturated and $X_i\cup I_i$ is a separate component; connectedness
and the distant geodesic vertices rule this out. Hence $q_i\geq1$,
while $x_i\geq1$ gives $q_i\leq k-1$.

For an internal index $i$, a vertex $v\in I_i$ cannot meet both
interfaces. Such neighbors, together with their neighbors in
$I_{i-1}$ and $I_{i+1}$, would give a path
\[
 z_{i-1},\ I_{i-1},\ F_{i-1},\ v,\ F_i,\ I_{i+1},\ z_{i+1}
\]
of length six between vertices at geodesic distance eight.
Every $v$ therefore chooses exactly one interface. Both choices
occur: if either side were empty, there would be no edge connecting
the vertices on opposite sides of that boundary in the chain of
types, contradicting connectedness. At the endpoints there is only
one available interface, so all their vertices use it.
\end{proof}

\begin{proposition}\label{primeobstruction}
If $k$ is prime and $r\geq2$, equality in Lemma~\ref{lower} is
impossible. For an odd prime $k$, every triangle-free $k$-regular
graph of diameter $4r$ therefore has at least $2k(r+1)+2$ vertices.
\end{proposition}
\begin{proof}
Put $b=|I_1^-|$ and $s=|F_0|$. Since $I_1$ is internal,
$1\leq b\leq k-1$. Count the edges incident with $F_0$. Its vertices
all have degree $k$, all neighbors lie in $I_0\cup I_1^-$, the former
contributes $kq_0$ edges, and the latter $bq_1$. Hence
\[
 ks=kq_0+bq_1.
\]
It follows that $k\mid bq_1$. For prime $k$, neither factor can be
divisible by $k$, because both lie between one and $k-1$. This is a
contradiction. Finally odd degree forces an even number of vertices
by the degree sum, so the strict lower bound increases by at least two.
\end{proof}

This argument uses only the first interface. The strict equality
obstruction is not an assertion that all the remaining interface data
are determined by that equation.

\section{A common metric lemma for the constructions}

We shall use a chain of $s=r-1\geq1$ nonempty independent sets
$X_1,\ldots,X_s$, called anchors in this section. There is a set $P_i$
of $k$ independent port vertices adjacent to every member of $X_i$.
Partition $P_i=P_i^-\sqcup P_i^+$ with both parts nonempty. For
$1\leq i<s$, an interface $U_i$ consists of independent vertices
whose neighbors lie in $P_i^+\cup P_{i+1}^-$, with every such port
meeting $U_i$. In the constructions each member of $U_i$ meets both
port sets, and paths between all choices of an adjacent pair of
anchors have length four. There are no other edges between the chain
pieces. A left endpiece meets the chain only through $P_1^-$ and
$X_1$; a right endpiece only through $P_s^+$ and $X_s$. Ports at those
ends belong to their respective endpieces.

\begin{lemma}\label{metric}
Suppose each vertex of either endpiece is within distance four of
every member of its adjacent anchor, and each endpiece has a specified
endpoint at distance exactly four from that anchor. Suppose also that
each interior interface vertex is within distance two of each adjacent
anchor. Then the full graph has diameter $4r$, provided it has no
additional edges beyond the described chain and endpieces.
\end{lemma}
\begin{proof}
Deleting $X_i$ separates everything strictly to its left from everything
strictly to its right: the two parts $P_i^-,P_i^+$ have no edge or
common neighboring piece except $X_i$. The distance between successive
anchors is exactly four. There is such a path through their port sets
and interface. There is no shorter path: their neighbor port sets
are disjoint, there is no edge between port vertices, and the chain
structure admits no bypass. Thus the two specified endpoints have
distance
\[
 4+4(s-1)+4=4r.
\]

For the upper bound, two vertices in the same endpiece are at distance
at most eight through any member of its anchor. Vertices in opposite
endpieces are at distance at most $4+4(s-1)+4=4r$.
Every remaining vertex is within two of an anchor: a port is within
one, an interface vertex within two, and a member of an anchor within
two of any selected representative of that anchor (two different
members share the port neighborhood). Two such remaining vertices are
therefore at distance at most $2+4(s-1)+2=4r-4$.
An endpiece vertex and a remaining vertex are at distance at most
$4+4(s-1)+2=4r-2$. These bounds include $s=1$, where there are no
interior interfaces. Since $r\geq2$, the same-endpiece bound eight
is at most $4r$. All pairs are covered, proving the claim.
\end{proof}

In both constructions the same side of the triangle-free incidence
structure contains all ports. The only possible extra edges within
an endpiece will be the explicitly described side edges in the
second construction. Those edges do not connect different pieces
of the chain, so the separator argument remains valid.

\section{Equality constructions for composite degree}

Let $k=pq$ with integers $p,q\geq2$. Put $s=r-1$. For each
$1\leq i\leq s$, take $|X_i|=k-q$ and $|P_i|=k$, with
\[
 |P_i^-|=p,\qquad |P_i^+|=k-p.
\]
Join every port of $P_i$ to every member of $X_i$. For $1\leq i<s$,
create $q$ interface vertices $U_i$, each adjacent to all vertices
of $P_i^+\cup P_{i+1}^-$. Their degrees are $k$, and every interior
port gains exactly $q$ neighbors from its interface.

\subsection*{The left endpiece}
Take an endpoint $e_0$, a set $A_0$ of $k$ row vertices and a side set
$T_0$ of $k$ vertices, labelled cyclically by $\mathbb Z/k\mathbb Z$.
Join $e_0$ to all rows. Join row $j$ to every side vertex except $j$;
this is the crown incidence $K_{k,k}$ minus its diagonal matching.
Partition $T_0$ into $p$ sets of size $q$, and join each port of
$P_1^-$ to one of those sets, using every set once. There are no
further left endpiece edges.

Each row has degree $1+(k-1)=k$, each side vertex has degree
$(k-1)+1=k$, and each port has $q+(k-q)=k$ after joining its anchor.
The endpoint also has degree $k$. Every side vertex is at distance
two from the anchor, every row at distance three, and the endpoint
at distance four. The exact endpoint distance follows from the only
permitted pattern $e_0$, row, side, port, anchor.

\subsection*{The right endpiece}
Take a set $E_r$ of $q-1$ endpoint twins, including a specified
$e_r$, a set $A_r$ of $k$ rows and a side set $T_r$ of $k$ vertices.
Join all of $E_r$ to all of $A_r$. Use any simple regular bipartite
incidence of degree $h=k-q+1$ between $A_r$ and $T_r$; for example,
on cyclic labels join row $j$ to side labels
$j,j+1,\ldots,j+h-1$ modulo $k$.

The $k-p=p(q-1)$ ports of $P_s^+$ form $q-1$ groups of $p$ ports.
For each group, partition $T_r$ into $p$ disjoint $q$-subsets and
join its ports to those subsets. Thus each port has $q$ side
neighbors and each side vertex has $q-1$ port neighbors. A row has
degree $(q-1)+h=k$ and a side vertex degree $h+(q-1)=k$.
Each endpoint twin has degree $k$, and each port again has degree
$(k-q)+q=k$ after its anchor is included.

All side vertices meet a port, so their distance to the anchor is
two. Since $h\geq1$, each row is at distance three and every endpoint
twin at distance four. In particular $e_r$ has distance exactly four.

\begin{proposition}\label{composite}
This construction is a simple connected $k$-regular graph of order
$2k(r+1)$, girth four and diameter $4r$.
\end{proposition}
\begin{proof}
All adjacencies are between distinct named vertex sets, without
repetition. The degree checks above also cover interior ports: all
have $k-q$ anchor neighbors and $q$ interface or endpiece neighbors.
An anchor vertex has all $k$ ports as neighbors. Interior interface
vertices have $(k-p)+p=k$ neighbors. The displayed incidences are
connected to the chain, so the whole graph is connected.

It is bipartite: one part contains the two row sets and all port
sets; the other contains anchors, interfaces, side sets and endpoint
twins. Moreover $k-q=q(p-1)\geq2$, so two members of any $X_i$
and two of its ports form a four-cycle. Thus the girth is four.

Count the first part: the two row sets contribute $2k$ and the
$s=r-1$ port sets contribute $(r-1)k$, for a total $(r+1)k$.
In the other part the endpoint twins contribute $1+(q-1)=q$,
the side sets $2k$, the anchors $(r-1)(k-q)$ and the interior
interfaces $(r-2)q$. Their sum is also $(r+1)k$.
This proves the claimed order. Successive anchors are linked by
length-four paths through any interface vertex; all interface
vertices are at distance two from either anchor. Both endpiece
radius and exact endpoint conditions were checked above. Apply
Lemma~\ref{metric} to obtain diameter $4r$.
\end{proof}

\section{The order-plus-two construction for every degree at least three}

Now let $k\geq3$ be arbitrary and again put $s=r-1$. Take
\[
 |X_i|=k-1,\quad |P_i^-|=k-2,\quad |P_i^+|=2.
\]
All ports of $P_i=P_i^-\sqcup P_i^+$ meet every member of $X_i$.
For each $1\leq i<s$, let $U_i$ be a single vertex adjacent to
all of $P_i^+\cup P_{i+1}^-$. Its degree is $2+(k-2)=k$;
each interior port gains its one required nonanchor neighbor.

\subsection*{The left endpiece}
Take one endpoint $e_0$, $k$ rows $A_0$ and a side set $T_0$ of
size $k$. Select distinct side vertices $x,y$ and add the side edge
$xy$. Join every one of the $k-2$ ports $P_1^-$ to $x$.
Join the endpoint to all rows. Each row has a $(k-1)$-subset of
$T_0$ as its side neighborhood, with the following omission multiset:
$x$ is omitted $k-1$ times and $y$ once. No other side vertex is
omitted. This specifies the incidence up to permuting the rows.

The side vertex $x$ meets one row, the side edge and $k-2$ ports,
giving degree $k$. The vertex $y$ meets $k-1$ rows and the side edge.
All other side vertices meet all $k$ rows. Rows have degree
$1+(k-1)=k$, the endpoint has degree $k$, and each port has
$1+(k-1)=k$ after joining its anchor.

There is no triangle. The only edge within the side set is $xy$,
and no row meets both its endpoints; no port meets $y$. All other
edges form an incidence graph with independent row and port sets.
For the radius check, ports are at distance one from the anchor,
$x$ at two and $y$ at three. The unique row omitting $y$ meets $x$
and is at distance three. Every other row meets $y$, so is within
four. Any other side vertex meets the first row and is within four.
The endpoint meets that row and has distance exactly four to the
anchor. Thus the whole endpiece is within four.

\subsection*{The right endpiece}
Take an endpoint $e_r$, $k$ rows $A_r$, and a side set of $k+1$
vertices. It consists of a star with center $w$ and $k-1$ leaves,
together with an isolated side vertex $f$. Choose a star leaf $y$.
The two ports of $P_s^+$ have sole side neighbors $w$ and $y$,
respectively. Join the endpoint to all rows. None of the rows meets
$w$. On the remaining $k$ side vertices (the leaves and $f$), each
row omits one vertex. The omission multiset is: $y$ twice, each other
leaf once, and $f$ never. Its total size is $2+(k-2)=k$, as required.

The center $w$ has degree $(k-1)+1=k$ from its star and port.
The leaf $y$ has $k-2$ row neighbors, one star neighbor and one
port neighbor. Every other leaf has $k-1$ row neighbors and its
star edge. The vertex $f$ meets all $k$ rows. Every row has
$k-1$ side neighbors and the endpoint, and each port has its side
neighbor and $k-1$ anchor neighbors. All degrees are therefore $k$.
The zero row column at $w$ is intentional and is allowed.

Again there is no triangle. A row cannot contain both endpoints
of a side edge because no row meets $w$, and a port has only one
side neighbor. For distances to the anchor, both ports are at one,
$w,y$ at two, and every other star leaf at three through $w$.
There is a row meeting $y$ since $k-2\geq1$; such a row is at
distance three. A row omitting $y$ contains another leaf, since
$k-2\geq1$, and is within four through that leaf and $w$.
The isolate $f$ meets every row and is within four through a row
meeting $y$. The endpoint is also at distance exactly four by
the same row. All right endpiece vertices are within four.

\begin{proposition}\label{plustwo}
For every $k\geq3$ and $r\geq2$, this construction is a simple
connected $k$-regular graph of order $2k(r+1)+2$, girth four and
diameter $4r$.
\end{proposition}
\begin{proof}
The incidences use distinct named vertex sets and contain no loops
or multiple edges. The endpiece degree checks above, the complete
anchor-port incidences, and the degree-$k$ singleton interfaces
verify every degree. Each endpiece is connected to its adjacent
anchor; the nonempty port sets and interfaces connect all anchors.

The only edges which could lie in a triangle outside the bipartite
chain are the left edge $xy$ and the right star edges. Their common
neighbor possibilities were excluded above, so the full graph is
triangle-free. Since $|X_i|=k-1\geq2$, two members of $X_i$ and
two ports form a four-cycle. Its girth is therefore exactly four.

The row and port sets contain $2k+(r-1)k=(r+1)k$ vertices. The
remaining vertices comprise two endpoints, the side sets of sizes
$k$ and $k+1$, the $(r-1)(k-1)$ anchor vertices and $r-2$ interface
vertices. Their total is
\[
 2+2k+1+(r-1)(k-1)+(r-2)=(r+1)k+2.
\]
Thus the full order is $2k(r+1)+2$. Each interface vertex is at
distance two from its adjacent anchors, which are linked by paths
of length four. Both endpieces meet the radius and exact endpoint
conditions proved above. Lemma~\ref{metric} gives diameter $4r$.
\end{proof}

\begin{proof}[Proof of Theorem~\ref{main}]
For composite $k$, Lemma~\ref{lower} and
Proposition~\ref{composite} give equal lower and upper bounds.
For prime $k\geq3$, the prime is odd, so
Proposition~\ref{primeobstruction} and
Proposition~\ref{plustwo} give equal lower and upper bounds.
Every upper construction has girth exactly four, as required.
\end{proof}

\section{Consequences and limits}

At diameter eight the theorem gives $n(3;4,8)=20$,
$n(4;4,8)=24$ and $n(5;4,8)=32$. These are deductions from
the common theorem, not imported numerical premises. The theorem
is a minimum-order result; it does not count all graphs attaining
the minimum. In particular it does not give the number of quintic
order-32 ordinary classes, the complete degree-three census, or
their source manuscript recoveries.

The equality data for $k=4,r=2$ do yield a small separate ordinary
classification: two eleven-vertex halves have two intrinsic
degree-two ports whose neighborhoods are two-subsets of a four-set;
their intersections have sizes $0,1,2$. The row omission multiset
is forced, and the full graph center is exactly the two anchor
twins. Gluing an unordered pair of these three halves gives six
ordinary classes. That full argument, including the converse and
metric checks, is retained separately in the diameter-eight note
\cite{eight}; it is not needed for the minimum theorem here.

No formal Lean certificate or computational control has been run
for these arguments. Any future computational use as a pruning
rule must cite an independently reviewed edition rather than
silently treating this draft as an accepted kernel theorem.

\begin{thebibliography}{9}
\bibitem{eight} Project hand proof, \emph{A regular-graph bound at
diameter eight and its degree-four equality case}, 13,975 bytes,
SHA-256 \texttt{26f46bb39cb0f22de00b3e7566096289ab18f645c90b3060130c14b3df6fa65e}.
Prepared independently from the coordinating agent's proposed hand
argument; independent review and adoption are separate.
\bibitem{five} Project retained fixed-geodesic normal forms for
quintic graphs of order 32, especially the independent-neighborhood
setup. Those detailed census proofs and original source attributions
are separate from the general-degree theorem of this note.
\end{thebibliography}
\end{document}
